\documentclass[10pt,a4paper]{amsart}
\usepackage[margin=19mm]{geometry}
\usepackage{amsmath,amssymb,mathtools}
\usepackage[scr=rsfs,cal=euler]{mathalfa}
\usepackage{xfrac}
\usepackage[unicode,hidelinks,bookmarks=false]{hyperref}
\newcommand{\ie}{i.e.\ }
\newcommand{\cf}{cf.\ }
\newcommand{\resp}{respectively\ }
\newcommand{\Subsection}{\subsection}
\providecommand{\nobreakdash}{\nobreak\hbox{-}}
\setlength{\emergencystretch}{2em}
\multlinegap0pt
\usepackage{amssymb}
\usepackage{mathtools}
\usepackage{bm}
\usepackage{xcolor}
\usepackage[shortlabels]{enumitem}
\newtheorem{proposition}{Proposition}
\newcommand{\Thetag}{\Theta_g}
\title{\textbf{On the encoding complexity of quantum numerical integration: an angle-structure characterization}\\[6pt]}
\author{F. Chinesta \and A. Falcó \and D. Falcó--Pomares}
\date{Source: arXiv:2604.24289v2 (CC BY 4.0)}
\begin{document}
\maketitle

\noindent\emph{Source and licence.} \textbf{On the encoding complexity of quantum numerical integration: an angle-structure characterization}\\[6pt]. Version arXiv:2604.24289v2 (CC BY 4.0), distributed by arXiv under the Creative Commons Attribution 4.0 International licence, http://creativecommons.org/licenses/by/4.0/, as recorded in the arXiv licence metadata for this exact version and shown on the abstract page https://arxiv.org/abs/2604.24289v2. Full licence notice: \texttt{licenses/arxiv-2604.24289-CC-BY-4.0.txt}.\par\smallskip

\noindent\textit{Editorial scope.} Counted unit: the article's proposition prop:encoding\_cost (source0.tex lines 858--879). The article's complete original proof of it is delivered with it (source0.tex lines 881--895, one \texttt{proof} environment of 15 lines). No mathematics is written, repaired or re-derived: the statement and the proof are the article's own lines, in the article's own notation, and no retained line is substituted or re-flowed.  Retained uncounted as the source-local context the proof needs: lines 611--618; lines 624--627.  Nothing was omitted from any retained span, so the package carries no omission record.  No retained line contains a citation, so the wrapper carries no bibliography and no bibliography entry.  No retained line contains a figure environment, an image-inclusion command or drawing code, so no image is delivered and the package declares no asset dependency.  Editorial formatting only, in addition: an AMS \texttt{amsart} preamble that restates the article's own declarations and macros used by the retained lines, namely the environments \texttt{proposition} on separate counters, together with the standard packages those lines need.  The retained environments print in this document as Proposition 1 in order of appearance, whereas the article prints the same items with the numbering of its own full document.  The complete native source of the article as distributed in the arXiv e-print for this exact version is bundled unchanged as \texttt{sources/199221/source0.tex}.
\begin{equation}
  G_g\,(\mathbf{u}_i\otimes\mathbf{u}_0^{(2)})
  = \mathbf{u}_i\otimes
    \bigl(\sqrt{1-g(x_i)}\,\mathbf{u}_0^{(2)}
         +\sqrt{g(x_i)}\,\mathbf{u}_1^{(2)}\bigr),
  \qquad i\in\{0,\ldots,N-1\},
  \label{eq:Gg_def}
\end{equation}

\begin{equation}
  A_g := G_g\,(H_n\otimes I_2) \;\in\; \mathrm{U}(2N),
  \label{eq:Ag_def}
\end{equation}

\begin{proposition}[Encoding cost equals angle-map complexity]
\label{prop:encoding_cost}
The encoding operator $G_g \in \mathrm{U}(2N)$ is block-diagonal with
respect to the decomposition $\mathbb{C}^{2N} = \bigoplus_{i=0}^{N-1}
\mathrm{span}\{\mathbf{u}_i\} \otimes \mathbb{C}^2$,
and the block at index~$i$ is the $2\times 2$ rotation matrix
\begin{equation}
  R_Y(\Thetag(b(i)))
  = \begin{pmatrix}
      \sqrt{1-g(x_i)} & -\sqrt{g(x_i)} \\
      \sqrt{g(x_i)}   & \sqrt{1-g(x_i)}
    \end{pmatrix},
  \label{eq:G_block}
\end{equation}
where $b(i)$ is the unique bit-string with $i(b(i))=i$.
Consequently, an efficient implementation of~$G_g$ is equivalent to
an efficient encoding of the map $b\mapsto\Thetag(b)$.
In particular, any upper bound on the circuit complexity of encoding
$\Thetag:\mathbb{B}^n\to[0,\pi]$ translates directly into an upper
bound on the gate complexity of~$A_g$; the algebraic structure of
$\Thetag$ therefore governs the encoding cost.
\end{proposition}

\begin{proof}
By \eqref{eq:Gg_def}, $G_g$ preserves each summand
$\mathrm{span}\{\mathbf{u}_i\}\otimes\mathbb{C}^2$, so it is
block-diagonal with a $2\times 2$ block at each index~$i$.
Writing $\theta_i := \Thetag(b(i)) = 2\arcsin(\sqrt{g(x_i)})$, the
half-angle identities give
$\cos(\theta_i/2) = \sqrt{1-g(x_i)}$ and
$\sin(\theta_i/2) = \sqrt{g(x_i)}$, so that block is exactly
$R_Y(\theta_i)$, which is~\eqref{eq:G_block}.
Since $A_g = G_g(H_n\otimes I_2)$ by~\eqref{eq:Ag_def} and
$H_n\otimes I_2$ costs $n$ gates independently of~$g$, every
factorisation of $b\mapsto\Thetag(b)$ into controlled rotations yields
a factorisation of $G_g$, and hence of $A_g$, with the same number of
non-trivial factors.
\end{proof}

\end{document}
